\documentclass[10pt]{article}
\usepackage[T1]{fontenc}
\usepackage[utf8]{inputenc}
\usepackage[french]{babel}
\usepackage{amsmath,amssymb,mathrsfs,mathpazo}
\usepackage{geometry}
\geometry{a4paper,margin=2cm}
\usepackage{graphicx,array,tabularx,tabularray,colortbl,multirow,multicol,booktabs,enumitem,pifont,ulem,fancybox,diagbox,cancel}
\usepackage{pstricks,pstricks-add,pst-plot,pst-tree,pst-node,pst-text,pst-func,pst-math,pst-bspline,pst-3dplot,pst-eucl,pst-solides3d}
\let\pscalculate\undefined
\usepackage{pst-bezier}
\usepackage{tikz,pgfplots,tkz-tab}
\usetikzlibrary{arrows,arrows.meta,patterns,calc,positioning,trees,decorations.markings,shapes.geometric,shapes.misc,shapes.arrows,intersections,decorations.pathmorphing,decorations.pathreplacing,backgrounds,shadows,babel}
\pgfplotsset{compat=1.18}
\usepgfplotslibrary{fillbetween}
\usepackage[np]{numprint}
\usepackage{listings,makecell,pas-tableur,twoopt}
\let\barre\undefined
\usepackage[pstricks]{tablvar}
\definecolor{darkgreen}{RGB}{0,100,0}
\providecommand{\red}{\color{red}}
\providecommand{\blue}{\color{blue}}
\providecommand{\green}{\color{green}}
\providecommand{\black}{\color{black}}
\providecommand{\white}{\color{white}}
\providecommand{\gray}{\color{gray}}
\providecommand{\cyan}{\color{cyan}}
\providecommand{\magenta}{\color{magenta}}
\providecommand{\yellow}{\color{yellow}}
\providecommand{\violet}{\color{violet}}
\providecommand{\orange}{\color{orange}}
\providecommand{\euro}{\texteuro}
\providecommand{\EUR}{\texteuro}
\providecommand{\ieme}{\textsuperscript{e}}
\providecommand{\degres}{\ensuremath{^\circ}}
\providecommand{\degree}{\ensuremath{^\circ}}
\providecommand{\vv}[1]{\overrightarrow{#1}}
\providecommand{\Coord}[1]{#1}
\providecommand{\decosix}{}
\providecommand{\starredbullet}{\ensuremath{\bullet}}
\pagestyle{empty}
\setlength{\parindent}{0pt}
\usepackage{adjustbox}\DeclareUnicodeCharacter{2236}{:}\DeclareUnicodeCharacter{2190}{\ensuremath{\leftarrow}}
\providecommand{\R}{\mathbb{R}}
\providecommand{\N}{\mathbb{N}}
\providecommand{\Z}{\mathbb{Z}}
\providecommand{\Q}{\mathbb{Q}}
\providecommand{\D}{\mathbb{D}}
\providecommand{\C}{\mathbb{C}}
\providecommand{\e}{{\mathrm{e}}}
\providecommand{\vect}[1]{\overrightarrow{#1}}
\providecommand{\vectt}[1]{\overrightarrow{\mathrm{#1}}}
\providecommand{\vvt}[1]{\overrightarrow{\mathrm{#1}}}
\providecommand{\np}[2][]{{#2~#1}}
\providecommand{\vv}[1]{{\overrightarrow{#1}}}
\renewcommand{\d}{{\mathrm{d}}}
\providecommand{\degres}{{^\circ}}
\providecommand{\degre}{{^\circ}}
\providecommand{\blue}{\color{blue}}
\providecommand{\red}{\color{red}}
\providecommand{\enskip}{\quad}
\providecommand{\llbracket}{\lbrack\!\lbrack}
\providecommand{\rrbracket}{\rbrack\!\rbrack}
\providecommand{\euro}{{\text{€}}}
\providecommand{\no}{{\text{n°}}}
\providecommand{\interval}[2]{{[#1;#2]}}
\providecommand{\intervFO}[2]{{[#1;#2[}}
\providecommand{\dotfill}{\ldots}

\providecommand{\ds}{\displaystyle}
\providecommand{\textding}[1]{#1}
\providecommand{\Oij}{\(\mathrm{O};\vec{\imath},\vec{\jmath}\)}
\providecommand{\Oijk}{\(\mathrm{O};\vec{\imath},\vec{\jmath},\vec{k}\)}
\providecommand{\Ouv}{\(\mathrm{O};\vec{u},\vec{v}\)}
\providecommand{\intervFF}[2]{\ensuremath{[#1;#2]}}
\renewcommand{\euro}{\ifmmode\text{\texteuro}\else\texteuro\fi}
\DeclareUnicodeCharacter{00B0}{\ifmmode^{\circ}\else\textdegree\fi}
\DeclareUnicodeCharacter{20AC}{\ifmmode\text{\texteuro}\else\texteuro\fi}
% Énoncés : Éducation nationale. Transcription et corrigés : APMEP, Denis Vergès et contributeurs. Découpage et classement : CoopMaths, Sébastien Lozano, Éric Elter et contributeurs. Adaptation de la présentation : Maths Explorer.
% CC BY-SA : https://coopmaths.fr/www/about/
% https://coopmaths.fr/alea/static/eam/2026/tex/eam_specifique_2026_09_antilles_3_cor.tex
\begin{document}
\noindent\textit{APMEP / CoopMaths. CC BY-SA. Présentation adaptée par Maths Explorer.}\par\medskip
\emph{Pour $t \in [0~;~+\infty[$, $f(t) = (2 - 0,5t)(7t+1)$ où $t$ est le temps en secondes et
$f(t)$ la hauteur du projectile en mètres.}

\begin{enumerate}
	\item
	\begin{enumerate}
		\item Pour tout réel $t$ de $[0~;~+\infty[$ :
		\[f(t) = (2 - 0,5t)(7t+1) = 2 \times 7t + 2 \times 1 - 0,5t \times 7t - 0,5t \times 1
		= 14t + 2 - 3,5t^2 - 0,5t\]
		donc $f(t) = -3,5t^2 + 13,5t + 2$.

		\item $f(0) = -3,5 \times 0^2 + 13,5 \times 0 + 2 = 2$.

Au moment de son lancement, le projectile se trouve à une hauteur de 2~mètres.
	\end{enumerate}
	\item
	\begin{enumerate}
		\item On résout graphiquement l'équation $f(t) = 12$ :
		
			\begin{center}
		\begin{tikzpicture}[x=3cm,y=0.35cm,>=stealth]
			
			\draw[gray!20,line width=0.3pt,xstep=0.1,ystep=1] (0,0) grid (4.3,16.5);
			\draw[gray!60,line width=0.5pt,xstep=0.5,ystep=5] (0,0) grid (4.3,16.5);
			
			\draw[->,line width=1pt] (0,0) -- (4.4,0) node[below right]{$t$};
			\draw[->,line width=1pt] (0,0) -- (0,17.2);
			
			\foreach \x in {0.5,1,1.5,2,2.5,3,3.5,4} \draw (\x,3pt) -- (\x,-3pt) node[below]{\footnotesize \np{\x}};
			\foreach \y in {5,10,15} \draw (2pt,\y) -- (-2pt,\y) node[left]{\footnotesize \np{\y}};
			\draw (0,0) node[below left]{\footnotesize 0};
			
			\draw[green!50!black,line width=1.2pt,domain=0:3,samples=100]
			plot (\x,{-3.5*(\x)^2+13.5*\x+2});
			\draw[green!50!black,line width=1.2pt,dotted,domain=3:4,samples=40]
			plot (\x,{-3.5*(\x)^2+13.5*\x+2});
			
			\draw[blue,line width=1pt] (0,12) -- (4,12) node[right]{$y=12$};
			
			\draw[blue,dashed] (1,12) -- (1,0);
			\draw[blue,dashed] (20/7,12) -- (20/7,0);
			\fill[blue] (1,12) circle(2pt);
			\fill[blue] (20/7,12) circle(2pt);
			\draw[blue] (1,0) node[below,yshift=-12pt]{$t_1=1$};
			\draw[blue] (20/7,0) node[below,yshift=-12pt]{$t_2\approx 2,85$};
		\end{tikzpicture}
	\end{center}
	
Les solutions sont
		\[t_1 = 1 \qquad\text{et}\qquad t_2  \approx 2,85.\]

La hauteur du projectile est supérieure ou égale à 12~m sur l'intervalle de temps $[1 ; 2,85]$ (en secondes).

		\item Le projectile touche le sol lorsque sa hauteur s'annule :
		\[f(t) = 0 \iff (2-0,5t)(7t+1) = 0 \iff 2-0,5t = 0 \text{ ou }7t+1=0.\]
		\[ \iff -0.5t = -2 \text{ ou } 7t = -1\]
		\[ \iff t = 4 \text{ ou } t = -\frac{1}{7}\]
		
		Seule la solution $t = 4$ appartient à l'intervalle $[0~;~+\infty[$.

		Le projectile touche le sol 4~secondes après son lancement.
	\end{enumerate}
\end{enumerate}

\end{document}
